% s-147-p147-m001; original PDF 147
\Aphorism{s-147-p147-m001}{232}
Der Autor, ein Mann, redet vom Weibe, und der Leser, ein Mann, ruft ihm begeistert zu: Ja, du kennst das Weib! Er sollte ihm ehrlicherweise nur sagen: dein Urtheil und Verständniss (oder Vorurtheil und Missverständniss) vom Weibe deckt sich mit dem meinen! Zwei Blinde einigen sich über eine Farbe, zwei Phantasten träumen das Gleiche: was ist damit für die Realität des Traumes, der Farbe bewiesen? — Daraus scheint zu folgen, dass man über das Weib — das Weib hören müsse, und nicht den Mann. Aber, so unwahrscheinlich es klingt, das Weib kann nicht reden! wo es redet, und von sich selbst redet, spricht es nur das Männergeschwätz nach. Bei allem Raffinement im Thun und Lassen gebricht es dem Weibe an jeglichem Raffinement im Darstellen; die männliche Optik hat bis zu dem Grade gesiegt, dass das Weib von sich selbst androcentrisch denkt. Es bleibt nichts übrig, man muss sich mit der männlichen \Foreign{fable convenue} hinsichtlich des Weibes begnügen; auch dieser Schiller hat die Schweiz nie gesehen, die er schildert, — aber in diesem Falle kennt sich der Schweizer selbst nur aus dem Wilhelm Tell. Nur in einem einzigen Gebiet, in dem eines specifisch weiblichen Cynismus, weiss es manche Baubo manchem Priapos zuvorzuthun: hier hat das Weib seinen eigenen Stil.
\par

